\subsection{Dedekind 整环的逼近定理}

在 Dedekind 整环中有一个“逼近定理”，它加强了第VI章第7节第2号定理1和第1节第5号命题9：

命题2. 设 A 为 Dedekind 整环，K 为其分式域，P 为 A 的非零素理想集合；对于 \(\mathfrak{p} \in \mathbf{P}\)，令 \(v_{\mathfrak{p}}\) 表示 A 的相应本性赋值。设 \(\mathfrak{p}_1, \ldots, \mathfrak{p}_q\) 是 P 中互不相同的元素，\(n_1, \ldots, n_q\) 是有理整数，\(x_1, \ldots, x_q\) 是 K 中的元素。则存在 \(x \in \mathbf{K}\)，使得 \(v_{\mathfrak{p}_i}(x - x_i) \geqslant n_i\)（\(1 \leqslant i \leqslant q\)），并且有 \(v_{\mathfrak{p}}(x) \geqslant 0\)，对 \(\mathfrak{p} \in \mathbf{P}\) 中除这些 \(\mathfrak{p}_i\) 之外的理想成立。

必要时把 \(n_{i}\) 换成更大的整数，可设它们全为正。先考察 \(x_{i}\) 属于 A 的情形；这显然等价于求一个 \(x \in A\)，满足同余关系式

\[
x \equiv x _ {i} (\mathrm{mod.} p _ {i} ^ {n _ {i}})
\]

于是 \(x\) 的存在性由第II章第1节第2号命题5得到。

现在转到一般情形。可写成 \(x_{i} = s^{-1}y_{i}\)，其中 \(s, y_{i}\) 属于 A；写成 \(x = s^{-1}y\) 后，问题等价于寻找一个 \(y \in A\)，使得一方面 \(v_{p_i}(y - y_i) \geqslant n_i + v_{p_i}(s)\)，另一方面，有 \(v_{p}(y) \geqslant v_{p}(s)\)，对 \(p \in P\) 中除 \(p_i\) 之外的理想成立；由于 \(v_{p}(s) = 0\)，除有限多个 \(p\) 外，故问题归结为上面的情形；命题得证。

命题2也可以解释为一个稠密性定理。确切地说，对所有 \(p \in P\)，令 \(\hat{K}_p\)（分别地，\(\hat{A}_p\)）表示 \(K\)（分别地，\(A\)）关于离散赋值 \(v_p\) 的完备化，并考察乘积 \(\prod_{p \in P} \hat{K}_p\)；这个乘积中的元素 \(x = (x_p)\) 称为 \(A\) 的限制阿黛尔，若 \(x_p \in \hat{A}_p\) 对所有 \(p \in P\)（除有限多个外）。显然，限制阿黛尔的集合 \(A\) 是 \(\prod_{p \in P} \hat{K}_p\) 的子环，其中包含乘积环 \(A_0 = \prod_{p \in P} \hat{A}_p\)。在 \(A_0\) 上考虑乘积拓扑；相对于该拓扑，\(A_0\) 是完备的；在 \(A\) 上存在唯一拓扑 \(T\)，它与 A 的加法群结构相容，并且使 \(A_0\) 中 0 的邻域构成 0 的基本邻域系 \(S\)。拓扑 \(T\) 与 \(A\) 的环结构相容；显然，一般拓扑第III章第6节第3号的公理（\(AV_{II}\)）成立，因为 \(T\) 在 \(A_0\) 上诱导的拓扑与 \(A_0\) 的环结构相容。另一方面，对于所有 \(x \in A\)，存在一个有限子集 \(J\)（属于 \(P\)），使得若记 \(J' = P - J\)，则 \(K_J \leqslant \prod_{p \in J} \hat{K}_p\)、\(A_{J'} = \prod_{p \in J'} \hat{A}_p\)，有 \(x \in K_{J} \times A_{J'}\)；由于 \(\hat{A}_{p}\) 对所有 p 在 \(\hat{K}_{p}\) 中是开集，\(\mathfrak{S}\) 是乘积拓扑在 \(K_{J} \times A_{J'}\) 上的 0 的基本邻域系。由于后一个乘积与其环结构相容，一般拓扑第III章所引位置的公理（AV \(_{I}\)）也成立，这就证明了断言。显然 \(A_{0}\) 是 A 的开子环，因此 A 也是完备环（一般拓扑第III章第3节第3号，命题4）。

对于所有 \(x \in \mathbf{K}\)，令 \(\Delta(x)\) 为元素 \((x_{\mathfrak{p}}) \in \bigcup_{\mathfrak{p} \in \mathbb{P}} \hat{\mathbf{K}}_{\mathfrak{p}}\)，满足 \(x_{\mathfrak{p}} = x\) 对所有 \(\mathfrak{p} \in \mathbb{P}\) 成立；由于除有限多个取值外 \(x_{\mathfrak{p}} \in \hat{\mathbf{A}}_{\mathfrak{p}}\)，故 \(\mathfrak{p}, \Delta(x) \in A\)；这样便定义了一个同态 \(\Delta: \mathbf{K} \to A\)，当 \(\mathbf{P} \neq \varnothing\) 时它是单射（即 \(A\) 不是域）；\(\Delta(\mathbf{K})\) 中的元素称为主限制阿黛尔，显然 \(\Delta(A) \subset A_0\)。

命题3. 环 \(A_0\)（分别地，\(A\)）可识别为 \(A\)（分别地，\(K\)）关于环拓扑的完备化；该拓扑的 0 的基本邻域系由所有满足 \(\neq (0)\) 的、属于 \(A\) 的整理想组成。

显然，在 A（或 K）上考虑的拓扑是豪斯多夫的。考虑第3号，关于 \(A_{0}\) 的断言由第III章第2节第13号命题17得到。因此，\(\Delta(A)\) 在 \(A_{0}\) 中稠密；同样，为证明 \(\Delta(K)\) 在 \(A\) 中稠密，注意对于所有 \(x = (x_{p}) \in A\)，\(p \in P\) 中满足 \(v_{p}(x_{p}) < 0\) 的 p 只有有限多个；由第1节第5号命题9，存在 \(s \in K\)，使 \(sx_{p} \in \hat{A}_{p}\) 对所有 \(p \in P\) 成立，换言之 \(\Delta(s)x \in A_{0}\)；由于乘以 \(\Delta(s)\) 是从 \(A\) 到自身的同态，只需利用 \(\Delta(A)\) 在 \(A_{0}\) 中稠密这一事实，便可推出 \(\Delta(K)\) 在 \(A\) 中稠密。

当然也可以利用命题2证明 \(\Delta(\mathbf{K})\) 在 A 中稠密。

现在考虑乘法群 \(\mathbf{SL}(n, A)\)，它由矩阵 \(U \in \mathbf{M}_n(A)\)（满足 \(\det(U) = 1\)）组成；若赋予 \(\mathbf{M}_n(A) = A^{n^2}\) 乘积拓扑，它在 \(\mathbf{SL}(n, A)\) 上诱导出与群结构 \(\mathbf{SL}(n, A)\) 相容的拓扑。只需验证映射 \(U \mapsto U^{-1}\) 在 \(\mathbf{SL}(n, A)\) 上连续；但由于 \(U\) 是幺模的，已知（《代数》第III章，第6节，第5号，公式(17)）\(U^{-1}\) 的元素是 \(U\) 的子式，因而是 \(U\) 的元素的多项式，这就证明了断言。若 \(K\) 被识别为 \(A\) 的子环，所用同构为 \(\Delta\)，则群 \(\mathbf{SL}(n, K)\) 是 \(\mathbf{SL}(n, A)\) 的子群。

命题4. 群 \(\mathbf{SL}(n, \mathbf{K})\) 在 \(\mathbf{SL}(n, A)\) 中稠密。

设 \(G\) 为 \(\mathbf{SL}(n, K)\) 在 \(\mathbf{SL}(n, A)\) 中的闭包；由于 \(K\) 在 \(A\) 中稠密（命题3），\(G\) 包含所有形如 \(I + a.E_{ij}\) 的矩阵，其中 \(i \neq j\) 且 \(a \in A\)。对于所有 \(p \in P\) 以及所有 \(\lambda \in \hat{K}_p\)，令 \(\lambda(p)\) 为限制阿黛尔 \(x = (x_q)_{q \in P}\)，满足 \(x_p = \lambda\) 且 \(x_q = 0\)（其中 \(q \neq p\)）；上述说明 \(G\) 包含所有形如 \(I + \lambda(p)E_{ij}\) 的矩阵，其中 \(i \neq j\)。但我们知道，形如 \(I + \lambda E_{ij}\) 的矩阵（其中 \(\lambda \in \hat{K}_p\)）生成群 \(\mathbf{SL}(n, \hat{\mathbf{K}}_p)\)（《代数》第III章）。对于每个矩阵 \(U \in \mathbf{SL}(n, A)\)，令 \(U_{p}\) 表示 U 在 \(\mathbf{SL}(n, \hat{\mathbf{K}}_{\mathfrak{p}})\) 中的典范像；因此可见，对所有 \(p \in P\)，G 包含矩阵 \(U \in \mathbf{SL}(n, A)\)，满足 \(U_{q} = I\) 对所有 \(q \neq p\)。由于 G 是群，它也包含所有矩阵 \(U \in \mathbf{SL}(n, A)\)，满足 \(U_{p} = I\)（除有限多个 \(p \in P\) 外）；现在，由 A 上拓扑的定义立即可知，这些矩阵的集合在 \(\mathbf{SL}(n, A)\) 中稠密。

